% STATUS: DRAFT
\chapter{Groups, unitary representations, and positive energy}\label{ch:groups}

\begin{physmotiv}
Symmetry and dynamics enter operator algebras as \emph{groups acting by automorphisms}. Time translation, boosts, gauge transformations and (crucially, later) the modular flow are all one-parameter groups. Two ideas from this chapter drive the rest of the book. First, a symmetry can be \emph{implemented} by unitaries on a given Hilbert space or it can fail to be (spontaneous symmetry breaking, \cref{ch:inequiv}). Second, a Hamiltonian that is \emph{bounded below} is a very strong constraint, and is the reason the de Sitter algebra will be of type II$_1$ rather than II$_\infty$ (\cref{ch:clpw}).
\end{physmotiv}

\section{Unitary representations}

A \emph{unitary representation} of a group $G$ on $\HH$ is a homomorphism $g\mapsto U(g)$ into the unitary group. For Lie groups and continuous symmetries we require \emph{strong continuity}: $g\mapsto U(g)\xi$ is norm-continuous for every $\xi$ (no stronger continuity can be demanded in infinite dimensions: the translation group on $L^2(\R)$ is not norm-continuous). By Stone's theorem (\cref{ch:spectral}) every one-parameter subgroup is generated by a self-adjoint operator: $U(t)=\ee^{-\ii tH}$.

\begin{center}
\renewcommand{\arraystretch}{1.25}
\begin{tabular}{@{}lll@{}}
\toprule
Group & Generator & Physics\\
\midrule
$\R$ (time) & $H$ & Hamiltonian\\
$\R^d$ (space) & $P_i$ & momentum\\
$SO(3)$ & $J_i$ & angular momentum\\
Poincar\'e & $P_\mu,M_{\mu\nu}$ & relativistic symmetry\\
$U(1)$ & $Q$ & charge ($\mathrm{spec}\,Q\subset\Z$)\\
modular flow & $K=-\log\Delta$ & ``hidden'' dynamics (\cref{ch:tomita})\\
\bottomrule
\end{tabular}
\end{center}

\subsection*{Haar measure}
To average over a group one needs an invariant measure. On $\R$, Lebesgue measure $\dd t$; on a compact group such as $U(1)$ or $SO(3)$, the normalized invariant measure; on a general locally compact group, the \emph{Haar measure}, unique up to a constant, which is left-invariant: $\int f(hg)\dd g=\int f(g)\dd g$ for all $h$. Groups for which left and right Haar measures agree are \emph{unimodular}; all groups in this book (abelian, compact, Poincar\'e) are unimodular except possibly the affine group $ax+b$, which appears in conformal/Möbius settings and as the group generated by modular flow and dilations. Group averaging is the physicist's ``sum over images'' and ``projection onto gauge-invariant states'': for a compact group,
\begin{equation}
P=\int_G\dd g\,U(g)
\end{equation}
is the projector onto the $G$-invariant vectors. For a non-compact group $P$ is not a bounded operator on $\HH$: see the exercise and \cref{ch:gravity_constraints}.

\section{Groups acting on algebras}

\begin{definition}
An \emph{action} of $G$ on a (\cstar\ or von Neumann) algebra $\M$ is a homomorphism $\alpha:G\to\Aut(\M)$, $g\mapsto\alpha_g$, where each $\alpha_g$ is a $^*$-automorphism (invertible, $^*$- and product-preserving). For a von Neumann algebra we also require $g\mapsto\alpha_g(a)$ to be continuous in the $\sigma$-weak topology.
\end{definition}

If $U(g)$ is a unitary representation on $\HH$ with $U(g)\M U(g)\ad=\M$, then $\alpha_g(a)=U(g)aU(g)\ad$ is an action; we say it is \emph{unitarily implemented} on $\HH$. An automorphism is \emph{inner} if it is implemented by a unitary \emph{in} $\M$ itself. The distinction between inner and outer is central:

\begin{center}
\emph{Inner:} the symmetry operator is an observable of the algebra (like $H$ for $\BH$).\\
\emph{Outer:} it acts on the algebra but its generator lies outside it.
\end{center}
In the algebra $\BH$ every automorphism is inner. The Rindler boost acts on the algebra of the right wedge $\A(R)$ as an outer automorphism: it is a symmetry of the algebra, but there is no operator in $\A(R)$ that generates it. This is the content of ``there is no density matrix in QFT'': for a density matrix $\rho$, the flow $\rho^{it}\,\cdot\,\rho^{-it}$ would be implemented by $\rho^{it}\in\M$.

\begin{whatwrong}
A symmetry of the algebra need not be implemented on a given representation. The automorphism group of the algebra may act in a representation $\pi_\omega$ by unitaries only when the state is invariant (more generally quasi-invariant) under it. If the ground state breaks the symmetry, no unitary $U(g)$ exists on $\HH_\omega$ and the symmetry is \emph{spontaneously broken}. This is the algebraic definition of SSB (\cref{ch:inequiv}), where it appears with infinite-volume ferromagnets as the main example.
\end{whatwrong}

\section{Positive energy}

In QM and QFT the Hamiltonian is required to be bounded below, $H\ge0$ (after subtracting the ground energy). In a Poincar\'e-covariant theory the stronger \emph{spectrum condition} says the joint spectrum of $P^\mu$ lies in the forward light cone. Important consequences:
\begin{itemize}
\item A bounded-below $H$ has a spectral projection $\chi_{[0,\infty)}(H)=\id$, so ``$\Tr\ee^{-\beta H}$'' makes sense for $\beta>0$ in good cases, and $\ee^{-sH}$ is a bounded operator for $s\ge0$ only. Analytic continuation of correlation functions into the lower half $t$-plane (the Euclidean regime) is available. (\emph{Contrast with the modular Hamiltonian $K=-\log\Delta$, which is \emph{unbounded in both directions}.})
\item A ground state exists as a lowest-energy vector (when $0$ is an eigenvalue of $H$), and thermal equilibrium states are characterized by stability properties that rely on energy being bounded below (\emph{passivity}, \cref{ch:kms}).
\end{itemize}

\subsection*{A preview of why the sign of $q$ matters}
We will study systems with a constraint of the form $H+q=0$ (or $H+q=\text{const}$), where $H$ is the Hamiltonian of a static patch (unbounded in both directions: it is the boost generator, equal to $K/\beta_{dS}$ up to normalization) and $q$ is the energy of an \emph{observer}, a clock with a physical Hamiltonian bounded below, $q\ge0$. Solving the constraint ties the value of $q$ to the spectrum of $H$, so that requiring $q\ge0$ restricts attention to a half-line of the spectrum of $H$ (signs and normalizations are fixed in \cref{ch:add_observer}). In algebraic language, the algebra of constraint-satisfying observables is a crossed product by the group $\R$ (\cref{ch:crossed}); without the condition $q\ge0$ the trace on the crossed product is infinite on the identity (type II$_\infty$); with $q\ge0$ it is finite (type II$_1$) (\cref{ch:clpw}). It is the same mathematics as the following elementary fact: $\int_{\R}\ee^{-x}\dd x=\infty$, whereas $\int_0^\infty\ee^{-x}\dd x=1$.

\begin{keyidea}
Symmetries act on algebras by automorphisms; they may be inner (generator in the algebra) or outer (not). Positive energy makes semigroups $\ee^{-sH}$ bounded, and will turn an infinite trace into a finite one.
\end{keyidea}

\section*{Exercises}
\begin{exercise}
Let $U(t)=\ee^{-\ii tH}$ with $H$ having purely continuous spectrum $\R$ (e.g.\ $H=p$ on $L^2(\R)$). Show that $\int_{-\infty}^\infty\dd t\,U(t)=2\pi\delta(H)$ as a distribution, i.e.\ it is not a bounded operator. For compact $G=U(1)$ and $H=Q$ with integer spectrum show instead that $\int_0^{2\pi}\frac{\dd\theta}{2\pi}\ee^{-\ii\theta Q}$ is the projector onto $Q=0$.
\end{exercise}
\begin{exercise}
Show that on $\BH=M_n(\C)$ every automorphism is inner (Skolem--Noether). Give an example of an automorphism of the algebra of an infinite spin chain (a \cstar-algebra) that is not inner, e.g.\ a lattice translation. Why can no unitary in the algebra implement it?
\end{exercise}
\begin{exercise}
Let $\M$ act on $\HH$ and $U(g)$ implement $\alpha_g$. Show that a different choice $U'(g)=U(g)z(g)$ with $z(g)\in\M'$ unitary also implements $\alpha_g$. When is this ambiguity absent? (Relevant to the freedom to redefine charges by gauge-algebra elements.)
\end{exercise}
\begin{exercise}
Compute $\int_0^\infty\ee^{-\beta\lambda}\dd\lambda$ and $\int_{-\infty}^{\infty}\ee^{-\beta\lambda}\dd\lambda$. Interpret the first as $\Tr\ee^{-\beta H}$ for a system with density of states 1 and spectrum $[0,\infty)$; explain why the second diverges and why this is the same issue as the type II$_\infty$/II$_1$ distinction above.
\end{exercise}
